% Préambule commun aux cours IPR Scholar.
\documentclass[11pt,a4paper,oneside]{article}
\usepackage[T1]{fontenc}
\usepackage[utf8]{inputenc}
\usepackage[french]{babel}
\usepackage{lmodern}
\usepackage{microtype}
\usepackage[a4paper,margin=2.35cm]{geometry}
\usepackage{amsmath,amssymb,mathtools}
\usepackage{booktabs,longtable,array}
\usepackage{xcolor}
\usepackage{hyperref}
\usepackage{tikz}
\usetikzlibrary{arrows.meta,positioning,shapes.geometric,calc}
\hypersetup{colorlinks=true,linkcolor=blue!55!black,citecolor=blue!55!black,urlcolor=blue!55!black}
\definecolor{iprblue}{HTML}{1F4E79}
\definecolor{iprgreen}{HTML}{2D6A4F}
\newcommand{\IPR}{\textsc{ipr}}
\newcommand{\coursmeta}[2]{\author{Projet \IPR{} -- documentation Scholar}\date{#1}\maketitle\begin{center}\small #2\end{center}\vspace{0.5cm}\tableofcontents\clearpage}
\newcommand{\resultat}[1]{\paragraph{À retenir.} #1}

\title{Moteur de rendu \IPR{} -- architecture, calcul et constantes}
\begin{document}
\coursmeta{22 octobre 2026}{Relecture sur \texttt{dev}, référence \texttt{a15f703}. Les campagnes citées gardent leur date et ne deviennent pas des tests de cette relecture.}

\section{Nature du moteur et frontière d'import}
Le moteur est un traitement d'image sur CPU, pas un moteur 3D. C++/Qt
orchestre des tampons OpenCV ; QML affiche les aperçus. LibRaw reste le
backend RAW courant, mais les adaptateurs enregistrés rejoignent la même
frontière que les lecteurs raster Qt : \texttt{ImageImport::Result}.
Un nouveau format exige un véritable lecteur, pas seulement une extension.

\begin{center}
\begin{tikzpicture}[node distance=8mm and 8mm,every node/.style={draw=iprblue,rounded corners,fill=blue!4,align=center,minimum height=8mm,font=\small}]
\node (raw) {RAW\\LibRaw / adaptateur};
\node[right=of raw] (import) {Import commun\\RGBA float Rec.2020};
\node[right=of import] (raster) {JPEG / PNG / TIFF\\Qt, orientation, ICC};
\node[below=of import] (render) {Proxy et retouches\\modules CPU};
\node[below=of render] (out) {Aperçu / export\\conversion sRGB};
\draw[-{Stealth}] (raw)--(import);
\draw[-{Stealth}] (raster)--(import);
\draw[-{Stealth}] (import)--(render);
\draw[-{Stealth}] (render)--(out);
\end{tikzpicture}
\end{center}

Le master est du RGBA float32 fini, Rec.2020 linéaire D65, alpha droit.
Les RGB négatifs et supérieurs à un sont permis à la frontière ; cela
n'assure pas que tous les opérateurs internes les préservent. Les poids
de luminance sont $Y=0{,}2627R+0{,}6780G+0{,}0593B$.
Les sorties sRGB utilisent leurs propres primaires et poids Rec.709.
Une étiquette ICC ne remplace jamais une conversion de valeurs~\cite{icc2022}.

\section{Développement caméra et profils}
Un RAW mosaïqué exige lecture, niveaux noir/blanc, WB et dématriçage~\cite{libraw}.
Le signal caméra pré-DCP, quand il existe, est immuable et distinct du master.
Le DCP applique matrices, gain de profil, tables et courbe dans leurs domaines
attendus, dont ProPhoto, puis rejoint Rec.2020. Sans DCP, le repli matriciel
du décodeur est converti vers le même master.

La calibration est demandée par défaut pour une nouvelle recette ; le
catalogue JSON schéma 4 doit reconnaître format, décodeur, boîtier, profil
et mode. Un appareil inconnu n'hérite pas des coefficients Canon. La case
avancée est un état demandé, pas une preuve d'application. DCP et LCP sont
des ressources distinctes ; leurs droits ne découlent pas de leur présence
dans un répertoire local.

\section{Ordre effectif des retouches}
\begin{enumerate}
\item Séparer RGB et alpha, sans modifier la source.
\item Avec contexte caméra DCP, rejouer WB manuelle et exposition avant le profil.
\item Appliquer les niveaux de gris, l'exposition sans contexte caméra, puis le contraste.
\item Sans contexte caméra, appliquer le repli température/nuance ; traiter les tons.
\item Appliquer saturation/vibrance, mélangeur huit bandes et grading.
\item Réduire le bruit, puis appliquer texture, clarté et voile.
\item Appliquer courbes luminance/RGB, flou, netteté et grain.
\item Réassembler l'alpha ; convertir et mapper la sortie dérivée vers sRGB.
\end{enumerate}
Cette liste suit \texttt{renderWorkingImage()}, pas une hiérarchie théorique
universelle. Le gamut mapping de sortie n'efface pas rétroactivement un
bornage intervenu dans une table ou une protection de traitement.

\section{Tons et domaines intermédiaires}
L'exposition applique $2^{EV}$ au signal linéaire approprié. Avec DCP, la
courbe appliquée ensuite peut modifier le rapport du TIFF : ne pas exiger
un doublement de ses codes à $+1$ EV. Le contraste courant est une S par
canal dans un domaine logarithmique, avec pivot de modèle $0{,}32$,
pas l'ancien pivot de travail $0{,}18$.

\texttt{ToneRecovery} emploie des courbes EV et les statistiques du neutre.
Les HL sont prédites à partir de scènes de référence voisines ; les courbes
sont échantillonnées puis rendues non décroissantes. Ombres, HL et blancs
négatifs agissent sur Y ; noirs et blancs positifs passent par les canaux
sRGB linéaires et reviennent au master. Ces intermédiaires ne constituent
pas un deuxième master. Les protections peuvent désaturer vers le blanc
ou renforcer la couleur lors d'un assombrissement. La réserve visuelle HL
du 21 octobre demeure ouverte ; le module local futur n'est pas livré.

\section{WB, couleur et détail}
Le repli RGB température/nuance construit une adaptation Bradford avec des
axes empiriques calibrés en sRGB, puis conjugue la matrice vers Rec.2020.
La normalisation Rec.709 de ces axes ne doit pas être remplacée arbitrairement
par les poids du master. La voie RAW avec contexte caméra adapte au contraire
le signal avant DCP. Les mêmes Kelvin ne garantissent pas les mêmes gains.

Le modèle global livré dans \texttt{dev} protège les faibles/fortes chromas
via un gain de vibrance dépendant de $(\max-\min)/\max$. Les variantes en
étude sur \texttt{patch\_color} ne sont pas présumées intégrées. Le mélangeur
HSV a huit bandes et une luminance sélective bornée à $\pm0{,}5$ EV.
Le grading colore trois zones tonales, sans détection d'objets.

Le débruitage sépare Y Rec.2020 et Co/Cg. Un MAD de détails Haar estime un
proxy par plan ; deux résidus gaussiens subissent un seuillage doux et une
protection de structure. Un estimateur de bruit n'est pas un détecteur
certain de pixels faux. Les blocs parfaitement constants sont ignorés
dans l'estimateur courant : c'est une limite à garder visible dans les études.
Les travaux statistiques non fusionnés de \texttt{patch\_noise} ne remplacent
pas ce modèle. La netteté est un masque flou adaptatif, pas un réseau IA.

\section{Registre des constantes vérifiables}
\begin{longtable}{p{0.20\linewidth}p{0.28\linewidth}p{0.41\linewidth}}
\toprule Valeur & Emplacement & Nature et effet \\\midrule\endhead
$0{,}2627$ ; $0{,}6780$ ; $0{,}0593$ & \texttt{WorkingColorSpace} & Poids de Y issus des primaires Rec.2020, pas gains de boîtier. \\
$65535$ & conversion entière 16 bits & Nombre maximal de codes ; ne borne pas tout tampon flottant à un. \\
$2000$--$50000$ K & WB / interface & Domaine accepté ; pas preuve de fidélité à tous les illuminants. \\
$4096$, $16$ EV & \texttt{ToneRecovery::Curve} & Résolution et étendue de la table ; coût et continuité à vérifier. \\
$0{,}80$, $0{,}995$ & épaule tonale & Protections empiriques de l'éclaircissement, pas limites physiques du capteur. \\
$0{,}05$, $0{,}5$ EV & mélangeur & Protection des gris et amplitude de luminance sélective. \\
$0{,}8/1{,}8$, $2{,}55/3{,}15$ & débruitage & Échelles gaussiennes et seuils Y/chroma empiriques. \\
$0{,}67448975$ & MAD Haar & Quantile d'une normale centrée réduite ; hypothèses de bruit à expliciter. \\
80 ms / 450 ms & rendu / sidecar & Temporisations d'événements, pas durées garanties des opérations. \\
80/75\,\% & contrôleur mémoire & Admission et hystérésis, pas plafond absolu imposé à l'OS. \\
64 Mio / 16 entrées & profils préparés & Budget partagé DCP/LCP, distinct du cache de pixels. \\
\bottomrule
\end{longtable}
Les constantes normatives, numériques, ergonomiques et empiriques n'ont pas
la même justification. Une valeur empirique doit garder son unité, son
domaine et la preuve qui motive son évolution.

\section{Asynchronisme et propriété mémoire}
Les travailleurs ne manipulent pas directement QML. Les révisions et jetons
d'annulation empêchent la publication d'un résultat périmé. Un proxy réduit
sert à l'interaction ; le zoom ne recrée pas des pixels pleine résolution.

\texttt{ImageCacheController} possède le LRU sur son thread, réserve la
mémoire des tâches et évince les entrées inactives les plus anciennes.
Le préchargement est de basse priorité et peut être refusé. L'ouverture
prioritaire peut attendre ; si IPR n'a plus rien à libérer, une admission sous
pression reste possible si le commit disponible le permet. Le seuil de
80\,\% n'est donc pas une garantie absolue. Les références actives peuvent
survivre à une éviction de l'entrée de cache.

Une entrée peut ne garder que proxy, signal caméra réduit et métadonnées.
L'export reconstruit alors la pleine résolution. Les profils préparés utilisent
un autre cache partagé borné, avec identité de contenu. Ni cache ne garantit
à lui seul une latence donnée.

\section{Persistance et limites}
Le sidecar courant v6 commence par \texttt{IPR1}, version, taille et date du
source. Il contient retouches, annotations et choix RAW/optique. Ces contrôles
d'identité ne sont pas une empreinte cryptographique des pixels source.
L'écriture différée passe par \texttt{QSaveFile}. L'INI schéma 2 conserve
uniquement la présentation et les préférences raccordées, indépendamment
des recettes photo.

Les JSON immuables des analyses font autorité sur les mesures ; les CSV,
courbes et textes sont dérivés. Un test de non-régression réussi ne démontre
ni une fidélité physique ni une parité Lightroom. Aucun export par lot,
affichage ICC par écran, masquage local ou traitement IA n'est déclaré
fonctionnel ici. Les chapitres réservés des manuels restent vides.

\resultat{Le master commun évite les moteurs parallèles par format. Des espaces intermédiaires spécialisés restent légitimes si leurs conversions, domaines et limites sont explicites.}
\bibliographystyle{plain}\bibliography{references}
\end{document}
